\subsection{発展的な読み物}

SWEBOK第01章で紹介される発展的な読み物を、日本語で読む目的に合わせて整理する。

\par\smallskip\noindent\refstepcounter{table}\label{tab:ch01-13}\textbf{表 \thetable: 第01章 ソフトウェア要求 表13}\par\smallskip
\begin{tabularx}{0.97\linewidth}{@{}p{.31\linewidth}X@{}}
\toprule
文献 & 読むと得られること \\
\midrule
\textit{BABOK Guide} & ビジネス分析全体の知識体系。ソフトウェア要求だけでなく、ビジネス課題、利害関係者、価値の分析を広く学べる。 \\
LaPlante and Kassab & 要求工学を体系的に学ぶための包括的な教科書。 \\
Robertson and Robertson & 要求プロセス、利害関係者、要求発見、要求記述を実務的に学べる。 \\
Gilb & 要求を測定可能にし、価値と品質を重視して記述する考え方を学べる。 \\
Wiegers, \textit{Software Development Pearls} & 長年の実務経験から得られた、要求を含むソフトウェア開発の実践的教訓を学べる。 \\
Fisher and Ury & 要求衝突の解決に役立つ交渉の基本を学べる。 \\
Ahmad & 電子的コミュニケーションを用いた面接で、暗黙知を引き出す方法に関する研究。 \\
\bottomrule
\end{tabularx}
